% JTDH manuscript file derived from the official English template.
% Compile with LuaLaTeX or XeLaTeX and biber.

\documentclass{article}
\usepackage{jtdh}
\usepackage[english]{babel}
\usepackage{url}
\usepackage[hang,flushmargin]{footmisc}
\usepackage{csquotes}
\usepackage{hyperref}
\hypersetup{colorlinks = true, linkcolor = black, urlcolor = gray, citecolor = black, anchorcolor = black}
\usepackage{booktabs}
\usepackage{amsfonts}
\usepackage{amssymb}
\usepackage{nicefrac}
\usepackage{xcolor}
\usepackage{graphicx}
\usepackage{tabularx}
\usepackage{pgfplots}
\pgfplotsset{compat=1.17}
\setlength{\floatsep}{8pt plus 2pt minus 2pt}
\setlength{\textfloatsep}{8pt plus 2pt minus 2pt}
\setlength{\intextsep}{8pt plus 2pt minus 2pt}
\usepackage[backend=biber, style=apa, uniquelist=false]{biblatex}
\addbibresource{references.bib}
\usepackage{textcomp}

\usepackage[]{gb4e}
\let\eachwordone=\sffamily
\let\eachwordtwo=\sffamily

\usepackage{enumitem}
\graphicspath{{img/}}

\aboverulesep=0ex
\belowrulesep=0ex

\title{A Community in Formation: Building and Analysing a Corpus of Slovene Posts on Bluesky}

\author{
  Nives HÜLL\textsuperscript{1}\orcidlink{0009-0004-1878-0102}
}
\institutions{
  \textsuperscript{1}University of Ljubljana, Ljubljana, Slovenia \rorlink{05njb9z20}\par
}

\begin{document}
\maketitle

\begin{abstract}
This paper presents a corpus of 141,013 public Slovene-language posts by 432 authors on Bluesky, spanning August 2023 to mid-April 2026. The corpus was built using a protocol-native ATProto workflow combining relay-based discovery of Slovene‑tagged posts, DID‑resolved backfill of surviving public post histories from authors’ Personal Data Servers, and post‑level language filtering validated through manual annotation.
As one of the first language-specific empirical studies of Bluesky use for a small European language community, the corpus documents the emergence of a new public space on an open‑protocol platform. Based on the collected data, a marked increase in posting activity is visible in late 2024. Across the collection period, reply posts account for a large share of the material (68.3\%), including in the months before this increase, indicating that conversational use was already present prior to the surge. Further observed characteristics include sparse and largely account‑concentrated hashtag use, and content sharing shaped primarily by Slovene news media, GIF‑based informal exchange, and Bluesky‑native tooling.
The paper illustrates how protocol‑level access in ATProto enables the construction of a high‑precision corpus for an under‑resourced language without relying on platform‑specific research APIs. At the same time, it raises legal and ethical questions that arise when small and highly identifiable language communities become reconstructable from public, decentralised infrastructures. The paper documents the collection workflow, its validation, and the basic temporal, interactional, and linking properties of Slovene‑language activity on Bluesky, offering a reproducible case study for digital humanities research on open‑protocol social media.
\end{abstract}
\keywordseng{Bluesky, ATProto, Slovene, corpus construction, social media}

\section{Introduction}

Bluesky and the wider ATProto ecosystem have become increasingly relevant as a source of public social-media data \parencite{failla2024bluesky,quelle2025bluesky}. For smaller languages such as Slovene, this opens up an opportunity to document an emerging online community on a new, decentralised platform, and to consider how such documentation should be done. The paper is also motivated by the author's long-standing interest in Slovene online communities, by familiarity with the emerging Slovene Bluesky space, and by a broader interest in what open-protocol social media makes visible and researchable.

This paper describes an attempt to build a Slovene Bluesky corpus from public data. The main challenges are not only in collection, but in deciding what should count as a Slovene post. Language tags are useful but imperfect, and posts by Slovene-linked users are not automatically Slovene. The approach combines protocol-aware collection with explicit validation, but it still requires drawing a practical line around what counts as a Slovene post.

Four questions guide the work. The first is whether a usable Slovene Bluesky corpus can be built from public ATProto data. The second concerns how far collection can move beyond a simple \texttt{langs=sl} strategy. The third asks which inclusion rules are reliable enough for a main corpus. The fourth asks what basic properties are visible in the resulting corpus.

This paper makes two contributions. First, it presents a protocol-native workflow for constructing a high-precision corpus of Slovene Bluesky posts from public ATProto data, combining tagged discovery, multi-PDS backfill, and validated post-level language filtering. Second, it uses the resulting corpus as an empirical case study of a small online community in formation. The analysis shows a sharp late-2024 growth surge, strongly conversational posting patterns, and a highly concentrated participation structure, together suggesting that the Slovene Bluesky sphere is emerging as a small conversational public space.

\section{Related Work}

\subsection{Slovene language resources and social-media corpora}

Written standard Slovene is well served by reference resources. Gigafida 2.0 \parencite{krek2020b-gf,gigafida-clarin} is the main reference corpus (1.2 billion tokens of newspaper, web, and book text, with automatic morphosyntactic annotation); ssj500k \parencite{krek2020a-ssj500k} provides gold-standard morphosyntactically annotated training data; siParl 3.0 \parencite{parliamentaryCorpus} covers parliamentary proceedings. The broader methodological context is surveyed in \textcite{studije2005korpusnem}.

The most directly relevant prior work is the JANES project (\textit{Jezikovna analiza nestandardne slovenščine}), which produced the first corpus of Slovene computer-mediated communication, drawing primarily on Twitter but also covering news comments and forum posts \parencite{erjavec2019janes}. JANES demonstrated both the feasibility of building Slovene social-media corpora and the specific challenges they pose: non-standard spelling, code-switching, short texts, and unreliable platform language tags. The present corpus continues in this tradition, shifting the target platform from Twitter to Bluesky and dealing with a community that has formed largely as a direct consequence of Twitter's decline.

\subsection{Bluesky and ATProto}

The migration of users from Twitter/X to alternative platforms accelerated sharply after Elon Musk's acquisition of Twitter in late 2022 and intensified further around the 2024 US presidential election. Bluesky became the main destination for a segment of this migration --- partly because its interface is close to what Twitter users were familiar with, making it considerably more accessible than federated alternatives such as Mastodon, which require users to choose a server instance and navigate inter-instance social friction. The corpus collected for this study directly captures the consequences of this shift, visible as a sharp late-2024 increase in Slovene post volume and a concentration of authors' first included posts in late 2024 (Section~\ref{sec:empirical}).

Technically, Bluesky is built on the AT Protocol (ATProto), a federated specification that separates identity from hosting: user identities are globally portable DIDs (decentralised identifiers) resolved through a public directory, and posts are stored on independently operated Personal Data Servers \parencite{atproto-spec}. From a corpus-collection standpoint, ATProto has one decisive advantage over proprietary platforms: all public content is accessible via a network-wide event firehose and per-account repository endpoints, without API quotas or paywalls. This makes collection reproducible at low cost, and the DID resolution mechanism allows a collector to follow an account regardless of which PDS hosts it --- a property exploited by the multi-PDS collection approach described in this paper.

Existing empirical work on Bluesky has so far operated at large scale and focused primarily on English-language content and network-level phenomena: \textcite{failla2024bluesky} assembled a dataset of 235 million posts from over four million users for platform-wide behavioural analysis, while \textcite{quelle2025bluesky} used Bluesky data to characterise network topology and political polarization. Language-specific corpus construction for a smaller community on ATProto has not previously been reported in the literature.

\subsection{Language identification in corpus construction}

Building a language-specific corpus from a multilingual stream typically combines two signals: platform-supplied language metadata and automatic language identification. Neither is reliable on its own. Platform metadata is user- or app-generated: in the present corpus, 41.4\% of posts carry both \texttt{sl} and \texttt{en} tags, almost entirely because Bluesky's default interface language is English, not because the posts are bilingual. Automatic language identification tools---langid.py \parencite{lui2012langid} and langdetect \parencite{nakatani2010langdetect} are used here---perform better at post level, but both struggle with short texts and with closely related languages; Slovene, Croatian, and Bosnian are routinely confused by standard detectors. Using consensus across multiple detectors and requiring agreement with author-supplied tags improves precision, which is the rationale for the tiered inclusion strategy described in Section~\ref{sec:corpus-construction}.

\section{Data Collection and Methodology}
\label{sec:corpus-construction}

The corpus was built through a layered pipeline: tagged discovery across the network, seed-author identification, full-history backfill from each author's home server, and post-level language filtering. Each phase is described below.

\subsection{Phase 1: Tagged discovery}

Two complementary collection methods were used to discover Slovene-linked posts across the Bluesky network.

First, a \textbf{historical crawl} enumerated all public repositories containing \texttt{app.bsky.feed.post} records via the ATProto relay endpoint \texttt{com.atproto.sync.listReposByCollection} on \texttt{bsky.network}. For each discovered repository, the author's DID was resolved to its home Personal Data Server (PDS), and the post records were fetched via \texttt{com.atproto.repo.listRecords} from that PDS. Only posts whose \texttt{langs} field contained the tag \texttt{sl} or any \texttt{sl-*} subtag were retained.

Second, a \textbf{live firehose} subscriber connected to the public repository event stream (\texttt{com.atproto.sync.subscribeRepos} on \texttt{wss://bsky.network}) and processed post creation, update, and deletion events in real time, again retaining only Slovene-tagged posts. The firehose subscriber also handled account-level events (deletions, deactivations, suspensions) by marking the affected posts accordingly in the local store.

Both sources stored their results in local SQLite databases, with posts keyed by their AT-URI and soft-deletion tracking for posts that were later removed by their authors.

\subsection{Phase 2: Seed author identification}

Every unique author DID that appeared in at least one Slovene-tagged post from either discovery source was collected into a seed author list. The resulting list defines the author population for the rest of the pipeline. The corpus is therefore not a sample of the Bluesky post stream, but a complete collection of posts from a defined set of discovered Slovene-linked accounts.

\subsection{Phase 3: Full-history backfill}

For each seed author, the complete surviving public post history was retrieved through ATProto repository access. Each author's DID was resolved through a public directory: \texttt{plc.directory} for \texttt{did:plc:} identifiers, or the \texttt{.well-known/did.json} endpoint for \texttt{did:web:} identifiers. The resolved PDS URL was cached locally to avoid repeated lookups. Posts were then fetched via paginated \texttt{com.atproto.repo.listRecords} calls to the author's home PDS.

Because resolution goes through the DID rather than through a fixed domain, the backfill was not restricted to accounts on \texttt{bsky.social}; accounts hosted on any public ATProto PDS were handled identically. This phase retrieved \textit{all} posts from each seed author, regardless of language tag, transforming the corpus from a tag-filtered post stream into an author-centred collection. Deleted posts that were no longer available at crawl time are not recoverable through this method and are out of scope.

The resulting expanded store contained 379,482 posts from 742 active authors, spanning the earliest surviving post (May 2023) through the collection date (April 2026).

\subsection{Phase 4: Post-level language filtering}

Each post in the expanded store was then classified individually to determine whether it should be included in the Slovene corpus. The classifier applied three binary signals:

\begin{itemize}
  \item \textbf{Tag signal}: does the post's \texttt{langs} field contain \texttt{sl} or \texttt{sl-*}?
  \item \textbf{langid signal}: does the langid.py language identifier \parencite{lui2012langid} assign the label \texttt{sl} to the post text?
  \item \textbf{langdetect signal}: does the langdetect library \parencite{nakatani2010langdetect} assign Slovene a probability of at least 0.90?
\end{itemize}

Before classification, URLs and @-mentions were stripped from the post text. Posts with fewer than 20 alphabetic characters after stripping were treated separately, because language identification on very short texts is unreliable. The langdetect library was seeded deterministically (seed~=~0) to ensure reproducible results.

For posts meeting the minimum length, the combination of signals determines the decision group. Table~\ref{tab:decision-groups} summarises the mapping.

\begin{table}[h]
\centering
\caption{Decision groups assigned by the language filtering step. Only groups marked as \textit{included} are present in the final corpus.}
\label{tab:decision-groups}
\begin{tabularx}{\linewidth}{lXXXl}
\toprule
Decision group & Tag & langid & langdetect & Status \\
\midrule
\texttt{core\_tag\_supported}        & \checkmark & \checkmark or --- & \checkmark or --- & Included \\
\texttt{review\_model\_consensus\_only} & --- & \checkmark & \checkmark & Included \\
\texttt{review\_langid\_only}         & --- & \checkmark & ---         & Included \\
\texttt{review\_langdetect\_only}     & --- & ---         & \checkmark & Included \\
\midrule
\texttt{review\_tag\_only}            & \checkmark & ---         & ---         & Excluded \\
\texttt{review\_short\_tagged}        & \checkmark & \multicolumn{2}{l}{(text too short)} & Excluded \\
\bottomrule
\end{tabularx}
\end{table}

The core group (\texttt{core\_tag\_supported}) contains posts where the author tagged the post as Slovene \textit{and} at least one language model agrees. The three model-supported review groups extend coverage to posts that were not tagged \texttt{sl} but are classified as Slovene by the models. Posts tagged \texttt{sl} with no model support (\texttt{review\_tag\_only}) and posts too short for reliable classification (\texttt{review\_short\_tagged}) were excluded after manual validation confirmed their higher noise rate (see Section~\ref{sec:validation}).

In addition to langid and langdetect, two further language identification tools --- fastText \parencite{joulin2017bag} and Lingua --- were evaluated on the manually validated samples. On the comparable subset, all four detectors showed the same observed precision; the relevant difference was recall. The langid--langdetect pair recovered more acceptable Slovene posts than the alternatives, so the original detector pair was retained for the final corpus build.

\subsection{Final corpus assembly}

The core and the three included review groups were merged into a single final corpus, deduplicated by post URI. Posts from multiple collection runs (historical crawl, live firehose, and incremental updates) were combined in this step, with the first occurrence of each URI retained. The final corpus contains 141,013 posts.

Each post is stored with its text, timestamp, author-supplied language tags, post-type flags (reply, quote), embed kind, and the per-post classifier outputs (decision group, langid label, langdetect probability). Author and post identifiers are pseudonymised in the public release. No linguistic annotation is included; the corpus is released as a raw text resource.

\subsection{Unit of analysis}

The unit of analysis is the individual post. Language classification is applied per post, not per author or per conversation thread. A post in English by a Slovene-linked author inside an otherwise Slovene thread is classified on its own text and excluded if the text is not Slovene. This means the corpus represents Slovene-language posting, not Slovene-linked authorship in general.

\section{Validation}
\label{sec:validation}

\subsection{Precision validation}

Validation was carried out through manual annotation of sampled posts. The strict core was evaluated using a random sample of 300 posts drawn uniformly from the \texttt{core\_tag\_supported} group. All 300 sampled posts were judged acceptable as Slovene, yielding an estimated precision of 100\% on the core layer.

The broader review layer was evaluated with a separate random sample of 300 posts. In that sample, 214 posts were labelled \texttt{Slovene-dominant}, 2 \texttt{Mixed-with-Slovene}, 40 \texttt{Not-Slovene}, and 44 \texttt{Undeterminable/too-short}, for an overall acceptability rate of 72\%. A breakdown by decision group revealed that the noise was concentrated in the tag-only and short-tagged subgroups: within this 300-post sample, every row drawn from the \texttt{review\_model\_consensus\_only}, \texttt{review\_langid\_only}, and \texttt{review\_langdetect\_only} groups was judged acceptable. Additional dedicated follow-up checks on the two smaller model-supported groups refined that picture: in the \texttt{review\_langid\_only} sample of 50 posts, 49 were labelled \texttt{Slovene-dominant} and 1 \texttt{Mixed-with-Slovene}; in the \texttt{review\_langdetect\_only} sample of 50 posts, 49 were labelled \texttt{Slovene-dominant} and 1 \texttt{Not-Slovene}.

These results support a final corpus that includes the strict core and all three model-supported review groups, while excluding the noisier tag-only and short-tagged material.

\subsection{False-negative estimation}

Precision validation confirms that posts included in the corpus are Slovene. A separate analysis estimated the false-negative rate, i.e.\ the fraction of Slovene posts from known authors that were excluded by the pipeline. A random sample of 200 posts from seed authors that (a) were not included in the final corpus, (b) carried no Slovene language tag, and (c) met the project's minimum-text rule of 20 alphabetic characters after URL and @-mention stripping was manually annotated. Of these, 17 posts (8.5\%) were judged to be Slovene-dominant or Mixed-with-Slovene.

Extrapolating to the full pool of 97,968 excluded untagged posts from known authors that met the same minimum-text rule suggests that roughly 8,300 Slovene posts from the same author set were not captured by the current validated rule. This corresponds to approximately 5.9\% of the size of the final corpus. The corpus should therefore be described as a high-precision collection of Slovene posts from discovered Slovene-linked authors under the current rule, not as a complete census of all Slovene writing on Bluesky.

\section{Empirical Analysis}
\label{sec:empirical}

\subsection{Corpus profile}

The final corpus contains 141,013 posts by 432 unique authors, spanning August 2023 to April 2026. Posts are distributed across four decision groups: 85,094 \texttt{core\_tag\_supported} (60.3\%), 46,134 \texttt{review\_model\_consensus\_only} (32.7\%), 5,632 \texttt{review\_langid\_only} (4.0\%), and 4,153 \texttt{review\_langdetect\_only} (2.9\%). The median post length is 86 characters.

\subsection{Platform adoption and growth}

Monthly post counts show one dominant structural break. In November 2024, Slovene post volume increased from 31 in October 2024 to 5,984 the following month, a 193-fold increase. Volume then stabilised at roughly 7,000--9,000 posts per month through most of 2025 and early 2026.

Author first-seen dates are likewise concentrated in late 2024: 172 authors have their first included post in November 2024 and 50 in December 2024. The corpus nevertheless contains low but non-zero pre-November 2024 activity, with 30 authors whose first included post falls between August 2023 and October 2024.

\subsection{Temporal patterns}

Posting activity follows a clear daily rhythm consistent with Central European Time (UTC+1/UTC+2). A morning peak occurs at 09:00--10:00~UTC (10--11~AM local) and a stronger evening peak at 18:00--19:00~UTC (7--8~PM local). Activity is lowest between 01:00 and 04:00~UTC. This pattern is consistent with a community of primarily Slovenian-timezone users posting during breaks and after work.

\subsection{Interactional structure}

Using a mutually exclusive hierarchy that counts reply-quote combinations as replies, the corpus is dominated by reply posts. Of the 141,013 posts, 96,278 (68.3\%) are replies, 41,517 (29.4\%) are original posts, and 3,218 (2.3\%) are non-reply quote posts. The median reply is 77 characters, compared to 114 characters for original posts. This reply-heavy pattern suggests that Slovene Bluesky use is strongly oriented towards dialogue rather than broadcasting.

The overall statistics mask a clear division between account types. Several accounts in the corpus---a national sports news outlet (\texttt{sportklubslovenija.bsky.social}), the SUP community account (\texttt{supdeska.bsky.social}), and an environmental organisation (\texttt{alpeadriagreen.bsky.social})---have reply rates of 0\%, functioning as pure broadcasters. Individual users, by contrast, often show reply rates exceeding 80--90\%. The inclusion of institutional broadcast accounts inflates the original-post count; among individual users the community is even more strongly conversational.

Posts with emoji occur at a rate of 19.5\%, with at least one external URL in the post text at 12.6\%, hashtags at 4.9\%, and \mbox{@-mentions} at 1.2\%. The low hashtag rate is consistent with the conversational character of the corpus. Across the whole corpus, 11.2\% of posts embed a structured external link card, 8.0\% include images, 2.3\% are non-reply quote posts embedding another post, 0.5\% contain video, and 0.2\% combine a quote-post with attached media. Video is notably rare.

\subsection{Author concentration and infrastructure}

Despite having 432 identified authors, posting activity is heavily concentrated. The Gini coefficient of posts per author is 0.82. The top 10\% of authors (43 users) account for 69.6\% of all posts, and just 22 authors are responsible for 50\% of the corpus. The median author contributed 35 posts over the entire collection period. This power-law distribution is characteristic of social media corpora generally, though the concentration here is particularly pronounced given the small and still-forming community.

Regarding platform infrastructure, 423 of 432 authors (97.9\%) are hosted on Bluesky's own PDS infrastructure (\texttt{*.host.bsky.network}), while 9 (2.1\%) are hosted on \texttt{eurosky.social}, a European-operated independent PDS. This confirms that the corpus is not restricted to a single server and that DID-based PDS resolution is operational in the collection workflow.

Handle choice is slightly more varied. Of the 432 authors, 395 (91.4\%) use a \texttt{*.bsky.social} handle, 30 (6.9\%) use a user-owned custom domain, and 7 (1.6\%) use a \texttt{*.eurosky.social} subdomain. The custom-domain group contributed 13,303 posts (9.4\% of the corpus). Two authors combine a custom domain handle with \texttt{eurosky.social} hosting.

Together, these figures show that most corpus authors use the default Bluesky configuration, while a smaller minority changed either their public handle, their hosting provider, or both.

\subsection{Topics: hashtags and linked domains}

Hashtag use is sparse but thematically informative. The most frequent hashtags are \texttt{\#supslovenija} (230 posts), \texttt{\#supdeska} (187), and \texttt{\#aquamarina} (177), all connected to the stand-up paddleboarding (SUP) community, which is evidently a well-represented subgroup. Other frequent hashtags include \texttt{\#pohvalanadan} (`praise a day'), \texttt{\#zelenjavajezakon} (`vegetables rule'), \texttt{\#catsofbluesky}, \texttt{\#slovenija}, \texttt{\#kava} (`coffee'), and \texttt{\#dobrojutro} (`good morning'), reflecting a mix of community ritual, current events, and platform culture.

The most linked external domain is \texttt{media.tenor.com} (2,613 posts), indicating frequent GIF sharing, followed by the news portal \texttt{n1info.si} (2,548), the Bluesky link-sharing tool \texttt{ebx.sh} (1,804), the Slovene tech news site \texttt{tehnozvezdje.si} (1,236), and YouTube/\texttt{youtu.be} (1,436 combined). The national public broadcaster \texttt{rtvslo.si} (702) also features prominently. GIF sharing as the single most common link type underscores the informal and conversational register of the corpus.

\subsection{Linguistic profile and multilinguality}

A notable feature of the corpus is the frequency of mixed language tagging. Of the 141,013 posts, 58,416 (41.4\%) were tagged by their authors with both \texttt{sl} and \texttt{en} language codes. By contrast, the language identification model langid detected a non-Slovene primary language in only 7.3\% of posts. The predominant non-Slovene labels assigned by langid were Croatian (\textit{hr}: 4,506 posts) and Bosnian (\textit{bs}: 2,285 posts)---closely related South Slavic languages for which cross-label confusion by automatic detectors is common. Genuine English-dominant content was detected in only 278 posts.

The high rate of \texttt{sl+en} author tagging without corresponding English detection by langid suggests a user-interface effect: Bluesky users with English-language app settings may have English added automatically to their post language tags, rather than this reflecting systematic code-switching at the post level. The corpus nevertheless contains posts with English lexical insertions within otherwise Slovene text, which are retained as authentic examples of the mixed-register writing characteristic of this community.

\section{Discussion}

The empirical findings in the previous section can be read in two ways: as descriptive statistics about a dataset, or as a portrait of a community caught in the process of forming. This section takes the second reading.

\subsection{A community with a visible history}

Most corpora are static snapshots --- the social processes that produced them are invisible in the data. This corpus is unusual in that its own growth structure narrates a social history. The near-zero activity before November 2024 and the sudden ×193 spike in that month are not noise: they mark the clearest structural break in the corpus. A majority of authors first appear in the corpus in November or December 2024, while lower-volume activity is already present from August 2023 onward. The corpus therefore functions simultaneously as a linguistic resource and as a record of a community in transition.

This has an implication for how the corpus should be used. Any NLP model trained on it, or any linguistic analysis conducted with it, is working with language produced under specific social conditions: by people who chose to move platforms, who were actively establishing a new online presence, and who were in many cases reconnecting with each other after a period of platform disruption. Whether these conditions leave traces in the language is an open question, but it is worth asking.

\subsection{Dialogue over broadcasting}

The most striking single finding is the reply rate: 68.3\% of all posts are replies. People are using the platform primarily to talk to each other, not only to publish at an audience.

This is reinforced by the broadcaster/conversationalist split. The small number of institutional accounts (sports news, environmental organisations, community groups) post exclusively original content. Among individual users the reply rate is often above 80--90\%. If institutional accounts are set aside, the community looks even more strongly conversational.

One interpretation is that this reflects the early-adopter character of the community: a small group of people who know each other, or quickly come to know each other, and who use the platform as a space for ongoing dialogue rather than mass communication. As the community grows, this dynamic may shift.

\subsection{Small but coherent}

With 432 authors and 141,013 posts, the Slovene Bluesky community is small by the standards of any major language platform. The Gini coefficient of 0.82 and the fact that 22 authors account for half the corpus confirm that a core group of highly active users sets the tone. The presence of a distinct SUP paddleboarding subgroup --- visible through hashtag patterns --- and the clear daily rhythm consistent with Slovenian working hours both suggest that, despite its small size, the community has genuine internal structure and is not simply a random collection of Slovene speakers who happen to be on the same platform.

\subsection{ATProto as a research infrastructure}

From a methodological standpoint, the experience of building this corpus confirms that ATProto's open design has real consequences for research. The absence of API rate limits or access agreements meant that a full backfill of the discovered Slovene-linked authors was achievable without institutional infrastructure. This stands in contrast to the conditions that prevailed on Twitter before the API changes of 2023, and even more so after. For researchers working on small or under-resourced languages --- where corpora are scarce and institutional resources are limited --- this openness matters. The Slovene community on Bluesky may be small, but it is unusually well-documented as a result.

\section{Limitations and Ethics}

\subsection{Coverage limits}

The corpus should not be described as a complete census of all Slovene Bluesky posts ever published. It covers currently recoverable public posts from discovered Slovene-linked accounts; deleted posts are excluded by design. The false-negative analysis (Section~\ref{sec:validation}) estimates an 8.5\% miss rate within the pool of excluded untagged posts from known authors that met the project's minimum-text rule; the true miss rate across the full user population is unknown and likely higher.

The methodology uses post-level rather than thread-level language classification. Slovene-linked users may post in English or other languages within otherwise Slovene conversations, and such replies are not automatically treated as Slovene corpus material.

\subsection{Open data and reproducibility}

One of the striking features of building this corpus is how open the data is. The AT Protocol was designed from the outset as a public infrastructure: there are no API keys required to access the public firehose, no rate limits on repository access, and no centralised gatekeeper deciding who can collect what. Anyone with a working internet connection and modest programming skills can reproduce the collection described here. For a small-language community like Slovene, this is an unusual and genuinely valuable property: it makes corpus construction feasible without institutional data-access agreements, and it allows the methodology to be fully transparent and reproducible.

This openness is what makes the research possible, and it is worth recognising as a design choice that distinguishes ATProto from its predecessors.

\subsection{Privacy and identifiability}

The same openness that makes collection easy raises a concern worth stating plainly. The Slovene Bluesky community, at the time of collection, consists of 432 identified authors. This is a small number. Someone reading this paper, or the corpus, could in most cases identify the individuals behind the accounts with little effort --- handles are publicly visible, many users link their real names or professional identities through custom domains, and the posting history of any individual is freely accessible through DID resolution. The corpus therefore represents not just a dataset but a detailed longitudinal record of specific, identifiable people.

This is not unique to Bluesky or to this corpus: it is a structural feature of any public social-media corpus at this scale. It is worth noting as a reminder that \textit{public} does not mean \textit{anonymous}. Users who post on Bluesky under their real names or with professionally linked handles are making a choice to be visible, but they may not have anticipated that their posting history would be compiled into a citable research corpus. The study follows standard practice by not reproducing full post texts in this article and recommending that any redistribution of the corpus take a conservative form, such as providing stable identifiers rather than full-text dumps, to limit the ease of re-identification for users who may later wish to reduce their public footprint.

Only public content was collected. Even so, reproducibility is better supported through code, decision rules, validation documentation, and aggregate statistics than through raw full-text redistribution.

\section{Future Work}\label{sec:future-work}

The most immediate next step is a legally and ethically defensible access format. The current public material consists of methodology and collection code only. A fuller corpus release would make the resource usable without re-running the entire pipeline, which matters for researchers without ATProto collection experience, but it should wait until the redistribution questions discussed in Section~\ref{sec:limitations} have been addressed. Possible formats include stable identifiers, strongly pseudonymised text, or controlled access rather than an unrestricted full-text dump.

A further step would be structured community consultation. Because this corpus captures a small and identifiable public, technically valid release formats may still differ in how acceptable they feel to the users whose posts are included. Future work could therefore combine legal and methodological analysis with lightweight community consultation on corpus access, quotation, and redistribution practices.

Adding linguistic annotation would extend the corpus considerably. A candidate pipeline has been prepared: CLASSLA for sentence segmentation, tokenisation, lemmatisation, and Slovene morphosyntactic tags \parencite{ljubesic2019neural}, with Trankit providing Universal Dependencies annotation \parencite{nguyen2021trankit}. This follows the broad direction established by the JANES project for Slovene user-generated content \parencite{erjavec2019janes}. The pipeline has not yet been run over the full corpus due to computational cost.

The corpus can also support follow-up work on non-standard Slovene processing, code-switching between Slovene and neighbouring languages, thread-level dialogue, and community formation during platform migration. The high reply rate (68.3\%) makes thread reconstruction especially relevant, while the temporal structure makes the corpus useful for longitudinal study.

Finally, the collection infrastructure can be extended longitudinally and applied to other small-language communities on Bluesky. That would make it possible to compare how under-resourced languages establish themselves on an open-protocol platform, and to test whether the Slovene pattern reported here is exceptional or part of a broader ATProto trajectory.

\section{Conclusion}

This paper has described a workflow for building a high-precision corpus of public Slovene Bluesky posts and used that corpus for an initial look at the interactional and temporal structure of Slovene activity on the platform.

The pipeline combines relay-based tagged discovery (historical crawl and live firehose), DID-based backfill from individual Personal Data Servers, and post-level language filtering validated through manual annotation. It is fully protocol-native: no special API keys, no institutional data-access agreements, and no restriction to a single hosting server. For small-language communities where research infrastructure is limited, this openness is a practical advantage.

The corpus of 141,013 posts by 432 authors shows a community whose recoverable activity expands sharply in late 2024, is dominated by dialogue rather than broadcasting, and is concentrated around a small core of highly active users. A clear Slovenian-timezone daily rhythm and a small set of users on independent infrastructure point towards a community with some internal structure, even if its small size makes finer-grained claims about subgroups difficult to support.

The corpus is best understood as a snapshot of currently recoverable public posts from discovered Slovene-linked authors, not a complete census of all Slovene writing on Bluesky. Deleted posts are excluded by design, and within the pool of excluded untagged posts from known authors that meet the project's minimum-text rule, the false-negative rate is estimated at approximately 8.5\%. The main next step is not only technical enrichment, but a redistribution model that makes the corpus useful without ignoring the rights and expectations of the people whose posts make it possible.


\section{Data and Code Availability}
The collection and anonymisation code is available at \url{https://github.com/hulln/slosky}. The corpus itself is not distributed as a full-text release at this stage; see Sections~\ref{sec:limitations} and~\ref{sec:future-work}. The corpus can be reconstructed from currently public ATProto data using the methodology and code described in this paper.

\acknowledgmenteng{
Generative AI tools were used selectively during the coding, implementation, and writing phases. All generated content was independently reviewed, revised where necessary, and finalised by the author.
}

\urlstyle{same}
\printbibliography

\newpage
\titleeng{Skupnost v nastajanju: gradnja in analiza korpusa slovenskih objav na omrežju Bluesky}
\begin{abstracteng}
Povzetek v slovenščini bo dodan v poznejši fazi priprave članka. Končna različica mora vključevati slovenski naslov, slovenski povzetek in slovenske ključne besede, skladno s predlogo JTDH.

Opomba k jeziku: Angleščina je v tem prispevku uporabljena izključno z namenom nagovarjanja širše mednarodne raziskovalne skupnosti na področjih digitalne humanistike, obdelave naravnega jezika in ATProta. Uporaba angleščine nikakor ne zmanjšuje pomena slovenščine v raziskovalnem in spletnem prostoru, katere vlogo še naprej dosledno in odločno zagovarjam.
\end{abstracteng}
\keywordsslo{Bluesky, ATProto, slovenščina, korpus, družbena omrežja}

\end{document}
